% Part: computability
% Chapter: computability-theory
% Section: rice-theorem

\documentclass[../../../include/open-logic-section]{subfiles}

\begin{document}

\olfileid{cmp}{thy}{rce}
\olsection{莱斯定理}

仔细想想就会发现，$\fn{Tot}$ 的具体细节并未在 \olref[tot]{prop:total} 的证明中起作用。
我们设计 $h(x,y)$，恰好在 $x \in K$ 时让它表现得像常函数 $j(y) = 0$；
但在该情况下，我们同样可以让它表现得像任何其他部分可计算函数。
这一观察使我们能够陈述一个更一般的定理，大致是说：可计算函数的任何非平凡性质都是不可判定的。

请记住，$\cfind{0}$、$\cfind{1}$、$\cfind{2}$、…… 是部分可计算函数的标准枚举。

\begin{thm}[莱斯定理]
  设 $C$ 是任意部分可计算函数的集合，并令 $A =
  \Setabs{n}{\cfind{n} \in C}$。如果 $A$ 是可判定的，那么 $C$ 要么是空集，要么是全体部分可计算函数的集合。
\end{thm}

{\em 指标集}是指满足如下性质的集合 $A$：如果 $n$ 和 $m$ 是“计算”同一函数的两个指标，那么 $n$ 和 $m$ 要么都属于 $A$，要么都不属于。不难看出，定理中的集合~$A$ 具有这一性质。反过来，如果 $A$ 是一个指标集，而 $C$ 是由这些指标计算的函数集合，那么 $A = \Setabs{n}{\cfind{n} \in C}$。

\begin{explain}
使用这一术语，莱斯定理等价于断言：任何非平凡的指标集都是不可判定的。
要理解该定理的含义，有必要强调一下\emph{程序}（比如你最喜欢的编程语言中的程序）与它们所计算的函数之间的区别。
程序（指标）作为句法对象，当然存在关于它们的可判定问题：这个程序的符号数超过150个吗？它超过22行吗？它有“while”语句吗？字符串“hello world”是否曾作为“print”语句的参数出现？
莱斯定理说的是，关于程序\emph{行为}的任何非平凡问题都是不可判定的。
这包括如下问题：该程序在输入为 $0$ 时停机吗？它是否会停机？它是否会输出偶数？
\end{explain}

\begin{proof}[莱斯定理的证明]
假设 $C$ 既不是空集，也不是全体部分可计算函数的集合，并令 $A$ 为 $C$ 中函数的指标集。我们将证明，如果 $A$ 是可判定的，我们就能解决停机问题；因此 $A$~不是可判定的。

不失一般性，我们可以假设处处无定义的函数 $f$ 不在 $C$ 中（否则，在下述论证中将 $C$ 替换为 $C$ 的补集）。设 $g$ 为 $C$ 中的任一函数。基本思路是，如果我们能判定~$A$，我们就能区分计算~$f$ 的指标与计算~$g$ 的指标；然后我们就可以利用这一点来解决停机问题。

具体做法如下。利用通用部分可计算函数，我们可以定义一个函数
\[
h(x,y) \simeq
\begin{cases}
\text{无定义} & \text{若 $\cfind{x}(x) \fundefined$} \\
g(y) & \text{否则。}
\end{cases}
\]
为了计算 $h$，首先尝试计算 $\cfind{x}(x)$；如果该计算停机，我们再计算~$g(y)$；如果{\em 该}计算也停机，我们就返回输出。更形式化地，我们可以写成
\[
h(x,y) \simeq \Proj{2}{0}(g(y),\fn{Un}(x,x)).
\]
其中 $\Proj{2}{0}(z_0, z_1) = z_0$ 是返回第 $0$ 个参数的二元投影函数，它是可计算的。

于是 $h$ 是部分可计算函数的复合，并且右边有定义且等于~$g(y)$ 当且仅当 $\fn{Un}(x,x)$ 和 $g(y)$ 都有定义。

注意，对于固定的~$x$，如果 $\cfind{x}(x)$ 无定义，那么 $h(x,y)$ 对每个~$y$ 都无定义；而如果 $\cfind{x}(x)$ 有定义，那么 $h(x,y) \simeq g(y)$。所以，对于任何固定的~$x$，$h(x,y)$ 要么表现得完全像 $f$，要么表现得完全像 $g$，而判定 $\cfind{x}(x)$ 是否有定义，就等价于判定这两种情况哪一种成立。但这又等价于判定 $h_x(y) \simeq h(x,y)$ 是否属于~$C$，而如果 $A$ 是可判定的，我们恰好可以做到这一点。

更形式化地说，由于 $h$ 是部分可计算的，存在某个指标~$e$，使得 $h$ 等于函数 $\cfind{e}$。根据 $s$-$m$-$n$ 定理，存在一个原始递归函数~$s$，使得对每个~$x$，$\cfind{s(e,x)}(y) = h_x(y)$。由此可得，对每个 $x$，如果 $\cfind{x}(x) \fdefined$，那么 $\cfind{s(e,x)}$ 与~$g$ 是同一个函数，因此 $s(e,x)$ 在 $A$ 中。另一方面，如果 $\cfind{x}(x) \uparrow$，那么 $\cfind{s(e,x)}$ 与~$f$ 是同一个函数，因此 $s(e,x)$ 不在~$A$ 中。换言之，对每个~$x$，$x \in K$ 当且仅当 $s(e,x) \in A$。如果 $A$ 是可判定的，那么 $K$~也是可判定的，这就产生矛盾。所以 $A$ 不是可判定的。
\end{proof}

莱斯定理非常强大。下面的直接推论展示了一些应用示例。

\begin{cor}
下列集合是不可判定的。
\begin{enumerate}
\item $\Setabs{x}{\text{$17$ is in the range of $\cfind{x}$}}$
\item $\Setabs{x}{\text{$\cfind{x}$ is constant}}$
\item $\Setabs{x}{\text{$\cfind{x}$ is total}}$
\item $\Setabs{x}{\text{whenever $y < y'$, $\cfind{x}(y) \fdefined$, and
    if $\cfind{x}(y') \fdefined$, then $\cfind{x}(y) < \cfind{x}(y')$}}$
\end{enumerate}
\end{cor}

\begin{proof}
它们都是非平凡的指标集。
\end{proof}

\end{document}